% Source: https://artofproblemsolving.com/community/c3h1309509_4_challenge_questions
% Stable identity: aops-topic-1309509-q04
Put $p=-x+y+z$, $q=x-y+z$, $r=x+y-z$. Factoring each difference of squares,
\[
x^2-(y-z)^2=(x-y+z)(x+y-z)=qr=a^2,\quad rp=b^2,\quad pq=c^2 .
\]
Multiplying, $(pqr)^2=(abc)^2$, so $pqr=\varepsilon\,|abc|$ with $\varepsilon=\pm1$; none of $p,q,r$ is $0$. Then
\[
p=\dfrac{pqr}{qr}=\dfrac{\varepsilon|abc|}{a^2},\qquad q=\dfrac{\varepsilon|abc|}{b^2},\qquad r=\dfrac{\varepsilon|abc|}{c^2}.
\]
Since $x=\dfrac{q+r}{2}$, $y=\dfrac{r+p}{2}$, $z=\dfrac{p+q}{2}$, the solutions are
\[
\boxed{(x,y,z)=\pm\dfrac{|abc|}{2}\left(\dfrac{1}{b^2}+\dfrac{1}{c^2},\ \dfrac{1}{c^2}+\dfrac{1}{a^2},\ \dfrac{1}{a^2}+\dfrac{1}{b^2}\right)},
\]
with the same sign in all three coordinates. Conversely these satisfy $qr=a^2$ etc., so both sign choices are solutions.
